
\documentclass[12pt,a4paper]{ctexart}

\usepackage[margin=2.3cm]{geometry}
\usepackage{amsmath}
\usepackage{array}
\usepackage{booktabs}
\usepackage{caption}
\usepackage{enumitem}
\usepackage{etoolbox}
\usepackage{fancyhdr}
\usepackage{float}
\usepackage{graphicx}
\usepackage{hyperref}
\usepackage{longtable}
\usepackage{tabularx}
\usepackage[table]{xcolor}
\usepackage{tikz}
\usetikzlibrary{arrows.meta,positioning,shapes.geometric,fit,calc}

\hypersetup{hidelinks}
\urlstyle{same}
\setlist{nosep}
\captionsetup{font=small}
\renewcommand{\arraystretch}{1.25}
\setlength{\parindent}{2em}
\setlength{\parskip}{0.3em}
\setlength{\tabcolsep}{4pt}
\emergencystretch=3em
\sloppy
\makeatletter
\g@addto@macro{\UrlBreaks}{\do\/\do\-\do\_\do\.\do\:}
\makeatother
\AtBeginEnvironment{longtable}{\small\setlength{\tabcolsep}{3pt}}
\AtBeginEnvironment{tabularx}{\small\setlength{\tabcolsep}{3pt}}
\AtBeginEnvironment{tabular}{\small\setlength{\tabcolsep}{3pt}}

\newcommand{\studentname}{许奕}
\newcommand{\studentid}{2351441}
\newcommand{\teamname}{许奕、王天一、马源胜、周由}
\newcommand{\coursename}{软件工程管理与经济}
\newcommand{\projectname}{基于大模型的建筑能源智能管理与运维优化系统}
\newcommand{\docversion}{V2.0}
\newcommand{\docdate}{2026年5月31日}
\newcommand{\code}[1]{{\footnotesize\nolinkurl{#1}}}
\newcolumntype{Y}{>{\raggedright\arraybackslash}X}
\newcolumntype{C}[1]{>{\centering\arraybackslash}p{#1}}
\newcolumntype{L}[1]{>{\raggedright\arraybackslash}p{#1}}

\tikzset{
  box/.style={rectangle, rounded corners, draw=cyan!55!black, fill=cyan!8, thick, minimum height=8mm, align=center, inner sep=5pt},
  process/.style={rectangle, rounded corners, draw=teal!55!black, fill=teal!8, thick, minimum height=8mm, align=center, inner sep=5pt},
  data/.style={cylinder, shape border rotate=90, aspect=0.25, draw=orange!70!black, fill=orange!10, thick, minimum height=12mm, align=center, inner sep=5pt},
  risk/.style={diamond, draw=red!60!black, fill=red!8, thick, aspect=2, align=center, inner sep=2pt},
  arrow/.style={-{Stealth[length=2.4mm]}, thick, draw=gray!70!black},
  note/.style={rectangle, draw=gray!60, fill=gray!6, rounded corners, align=left, inner sep=6pt}
}

\pagestyle{fancy}
\fancyhf{}
\fancyhead[L]{\coursename}
\fancyhead[R]{\reporttitle}
\fancyfoot[C]{\thepage}
\setlength{\headheight}{15pt}

\title{
    \vspace{-2em}
    \heiti \zihao{2} \reporttitle \\
    \vspace{1em}
    \zihao{4} 课程期末项目正式交付文档
}
\author{
    \songti \zihao{4}
    项目名称：\projectname \\
    项目组成员：\teamname \\
    负责人：\studentname \quad 学号：\studentid
}
\date{\docdate}

\newcommand{\reporttitle}{SRS 软件需求规格说明书}
\begin{document}

\maketitle
\thispagestyle{fancy}

\begin{abstract}
本文档为课程期末项目《\projectname》的软件需求规格说明书，采用正式 SRS 结构描述系统目标、业务背景、用户角色、系统边界、功能需求、非功能需求、数据需求、外部接口和验收标准。本文档的需求基线来自项目最终实现版本：系统已形成 FastAPI 后端、Vue 前端、建筑能耗样例数据、异常分析、三维楼层态势、工单闭环、知识库问答、可选外部大模型和 MCP Server。文档中的功能项与代码仓库、测试脚本和最终演示路径保持一致，可作为课程评审、系统验收和后续维护的需求依据。
\end{abstract}

\tableofcontents
\newpage

\section{文档控制}
\begin{longtable}{p{4cm}p{10.5cm}}
\toprule
项目 & 内容 \\
\midrule
文档名称 & SRS 软件需求规格说明书 \\
文档版本 & \docversion \\
适用阶段 & 期末项目最终提交、系统演示、验收评审、后续维护 \\
项目范围 & Web 前端、FastAPI 后端、MCP Server、样例数据集、知识库、测试脚本和交付文档 \\
关联文档 & SDD 软件设计说明书、SEE 软件经济分析与评价、SEM 软件工程管理说明、测试与验收报告、用户手册与部署说明 \\
参考资料 & IEEE 风格 SRS 文档结构；同类课程项目文档组织方式参考 \url{https://github.com/L-Dramatic/Intelligent-Contract-Management-System} \\
\bottomrule
\end{longtable}

\section{引言}
\subsection{编写目的}
本文档用于明确本系统的业务问题、用户需求、软件功能、外部接口、数据约束、非功能要求和验收标准。它解决三个问题：第一，系统到底要完成哪些业务任务；第二，哪些功能属于本次课程项目边界，哪些属于后续扩展；第三，评审人员如何根据文档检查系统是否已经满足课程项目要求。

\subsection{项目背景}
校园建筑能源管理通常存在数据分散、异常发现滞后、设备状态难以追踪、人工巡检依赖经验、节能措施缺少量化依据等问题。传统的表格统计方式能够记录用电、用水和设备运行数据，但很难支撑快速定位“哪栋楼、哪一层、哪类设备、什么时段出现异常”。本项目以校园建筑能源管理为场景，构建一个可运行的智能管理与运维优化平台，将数据查询、统计分析、异常解释、三维风险态势、工单处理、决策报告和智能问答组合成完整业务闭环。

\subsection{系统目标}
系统目标不是单纯展示几张图表，而是完成从数据到处置的闭环。具体目标如下：
\begin{enumerate}
    \item 建立标准化建筑能耗样例数据集，覆盖多建筑、多时段、多设备和多指标。
    \item 提供按建筑、楼层、时间范围的查询和导出能力，支撑数据追溯。
    \item 提供总览、趋势、建筑对比、COP 排名、异常原因、楼层分析和设备分析。
    \item 使用可解释规则定位异常，并输出触发规则、指标证据和处理建议。
    \item 使用三维楼宇视图展示楼层健康状态，并与统计分析页面联动。
    \item 建立异常工单闭环，使异常能够被分派、跟踪和完成。
    \item 提供本地知识库问答和可选外部大模型增强，提升运维解释能力。
    \item 提供 MCP Server，使 AI 客户端能够通过协议调用项目数据和分析能力。
\end{enumerate}

\section{产品范围与业务边界}
\subsection{产品定位}
本系统定位为课程项目级“建筑能源智能管理与运维优化原型系统”。它不直接接入真实楼宇自控系统，也不承担法定能源审计责任；它面向教学演示、工程管理课程评审和后续原型扩展，重点体现软件工程中的需求、设计、实现、测试、部署、经济分析和项目管理闭环。

\subsection{业务边界}
\begin{table}[H]
\centering
\caption{系统业务边界}
\begin{tabularx}{\textwidth}{p{3.2cm}YY}
\toprule
类别 & 纳入本期范围 & 不纳入本期范围 \\
\midrule
数据来源 & 样例 CSV 数据、数据字典、知识库 Markdown、运行期工单 JSON & 实时电表、真实楼宇自控系统、真实收费账单 \\
业务功能 & 查询、统计、异常、工单、报告、问答、MCP 接入 & 多租户权限、组织审批流、企业级审计 \\
算法能力 & 可解释规则、COP 计算、建筑基线、楼层聚合 & 大规模模型训练、在线学习、真实预测竞赛模型 \\
部署形态 & 本地开发部署、课程演示部署、MCP 本地接入 & 高可用生产集群、移动端 App、IoT 边缘网关 \\
\bottomrule
\end{tabularx}
\end{table}

\subsection{系统上下文图}
\begin{figure}[H]
\centering
\resizebox{0.96\textwidth}{!}{%
\begin{tikzpicture}[node distance=12mm and 16mm]
\node[box, minimum width=29mm] (manager) {能源管理人员};
\node[box, below=of manager, minimum width=29mm] (worker) {运维人员};
\node[box, below=of worker, minimum width=29mm] (teacher) {评审人员};
\node[process, right=32mm of worker, minimum width=42mm, minimum height=24mm] (system) {建筑能源智能\\管理平台};
\node[data, right=30mm of system, minimum width=32mm] (dataset) {能耗 CSV\\工单 JSON};
\node[data, above=of dataset, minimum width=32mm] (kb) {知识库\\Markdown};
\node[box, below=of dataset, minimum width=32mm] (llm) {外部大模型\\可选};
\node[box, below=of system, minimum width=42mm] (mcp) {支持 MCP 的\\AI 客户端};
\draw[arrow] (manager) -- node[above, sloped]{查看报表} (system);
\draw[arrow] (worker) -- node[above, sloped]{处理异常} (system);
\draw[arrow] (teacher) -- node[above, sloped]{验收演示} (system);
\draw[arrow] (system) -- node[above]{读写} (dataset);
\draw[arrow] (system) -- node[right]{检索} (kb);
\draw[arrow] (system) -- node[right]{增强问答} (llm);
\draw[arrow] (mcp) -- node[right]{协议调用} (system);
\end{tikzpicture}
}
\caption{系统上下文与外部角色}
\end{figure}

\section{用户角色与使用场景}
\begin{longtable}{p{3.2cm}p{4.4cm}p{6.8cm}}
\caption{用户角色定义}\\
\toprule
角色 & 核心诉求 & 典型任务 \\
\midrule
\endfirsthead
\toprule
角色 & 核心诉求 & 典型任务 \\
\midrule
\endhead
能源管理人员 & 掌握整体能耗、异常趋势和节能空间 & 查看总览 KPI、建筑对比、COP 排名、运营报告和节能建议 \\
运维人员 & 快速定位异常并形成处理记录 & 查看异常明细、解释异常、创建工单、更新工单状态、确认处理结果 \\
项目管理者 & 评估系统价值和交付完整性 & 查看经济分析、工程管理、测试报告和最终交付材料 \\
任课教师与评审人员 & 验证系统是否满足期末项目要求 & 按 README 运行系统，核验 SRS、SDD、SEE、SEM、源码和系统展示路径 \\
AI 客户端或智能体 & 通过标准协议调用项目能力 & 使用 MCP Tools 查询数据、获取异常解释、生成运营报告和检索知识库 \\
\bottomrule
\end{longtable}

\section{业务流程需求}
\subsection{核心业务流程}
\begin{figure}[H]
\centering
\resizebox{0.98\textwidth}{!}{%
\begin{tikzpicture}[x=1cm,y=1cm,every node/.style={font=\small}]
\node[process, minimum width=2.4cm, minimum height=0.9cm] (data) at (0,0) {1 数据读取\\字段校验};
\node[process, minimum width=2.7cm, minimum height=0.9cm] (query) at (3.0,0) {2 查询筛选\\建筑/楼层/时间};
\node[process, minimum width=2.8cm, minimum height=0.9cm] (analysis) at (6.2,0) {3 统计分析\\趋势/COP/对比};
\node[risk, minimum width=2.3cm, minimum height=1.0cm] (risk) at (9.4,0) {4 异常\\判断};
\node[process, minimum width=2.7cm, minimum height=0.9cm] (explain) at (9.4,-2.1) {5 异常解释\\证据与建议};
\node[process, minimum width=2.7cm, minimum height=0.9cm] (order) at (6.2,-2.1) {6 生成工单\\分派负责人};
\node[process, minimum width=2.7cm, minimum height=0.9cm] (finish) at (3.0,-2.1) {7 处理完成\\状态持久化};
\node[process, minimum width=2.8cm, minimum height=0.9cm] (report) at (0,-2.1) {8 运营报告\\管理建议};

\draw[arrow] (data) -- (query);
\draw[arrow] (query) -- (analysis);
\draw[arrow] (analysis) -- (risk);
\draw[arrow] (risk) -- node[right, fill=white, inner sep=1pt]{异常} (explain);
\draw[arrow] (explain) -- (order);
\draw[arrow] (order) -- (finish);
\draw[arrow] (finish) -- (report);
\draw[arrow, dashed] (analysis.south) -- node[right, fill=white, inner sep=1pt]{统计结果} (order.north);
\end{tikzpicture}
}
\caption{从数据到工单的业务闭环}
\end{figure}

\subsection{用例列表}
\begin{longtable}{p{2cm}p{3.2cm}p{3cm}p{5.8cm}}
\caption{主要用例}\\
\toprule
编号 & 用例名称 & 主要参与者 & 成功结果 \\
\midrule
\endfirsthead
\toprule
编号 & 用例名称 & 主要参与者 & 成功结果 \\
\midrule
\endhead
UC-01 & 查看系统总览 & 能源管理人员 & 展示记录数、建筑数、平均 COP、异常数量和总能耗 \\
UC-02 & 浏览能耗记录 & 能源管理人员 & 按建筑、楼层、时间和数量筛选记录，表格字段完整 \\
UC-03 & 导出查询数据 & 能源管理人员 & 下载当前筛选范围 CSV，文件可被表格软件打开 \\
UC-04 & 查看统计分析 & 管理人员 & 展示趋势、异常、楼层、设备、推荐和全局对比 \\
UC-05 & 查看健康楼层 & 管理人员 & 绿色楼层只显示健康结论，不出现无意义空模块 \\
UC-06 & 点击三维楼层 & 运维人员 & 跳转统计分析并自动锁定对应建筑和楼层 \\
UC-07 & 解释异常记录 & 运维人员 & 输出触发规则、关键指标、结论和建议动作 \\
UC-08 & 创建异常工单 & 运维人员 & 选择异常、分配负责人，生成处理中工单 \\
UC-09 & 完成异常工单 & 运维人员 & 工单变为已完成并持久化，刷新后仍保留 \\
UC-10 & 生成决策报告 & 管理人员 & 输出风险摘要、工单状态、优化建议和收益测算 \\
UC-11 & 使用智能问答 & 管理人员 & 基于知识库和项目数据回答能源运维问题 \\
UC-12 & MCP 查询项目能力 & AI 客户端 & 通过 MCP Tools 获取数据、统计、异常解释和报告 \\
\bottomrule
\end{longtable}

\subsection{用例图}
\begin{figure}[H]
\centering
\resizebox{0.98\textwidth}{!}{%
\begin{tikzpicture}[every node/.style={font=\scriptsize}]
\tikzstyle{actorbox}=[draw=teal!75!black, rounded corners=3pt, fill=teal!6,
    minimum width=23mm, minimum height=8mm, align=center]
\tikzstyle{usecase}=[ellipse, draw=teal!70!black, fill=teal!7,
    minimum width=27mm, minimum height=8mm, align=center]
\tikzstyle{riskcase}=[ellipse, draw=red!65!black, fill=red!7,
    minimum width=27mm, minimum height=8mm, align=center]
\tikzstyle{workcase}=[ellipse, draw=orange!80!black, fill=orange!9,
    minimum width=27mm, minimum height=8mm, align=center]
\tikzstyle{aicase}=[ellipse, draw=cyan!70!black, fill=cyan!7,
    minimum width=27mm, minimum height=8mm, align=center]
\tikzstyle{verifycase}=[ellipse, draw=gray!70, fill=gray!8,
    minimum width=27mm, minimum height=8mm, align=center]

\draw[rounded corners=6pt, draw=gray!60, fill=gray!3] (0,-4.95) rectangle (10.5,2.55);
\node[font=\small] at (5.25,2.25) {建筑能源智能管理平台};

\node[actorbox] (energy) at (-1.65,0.75) {能源管理\\人员};
\node[actorbox] (reviewer) at (-1.65,-3.7) {评审人员};
\node[actorbox] (operator) at (12.15,-0.8) {运维人员};
\node[actorbox] (agent) at (12.15,-3.3) {AI 客户端};

\node[usecase] (uc01) at (2.0,1.35) {查看总览};
\node[usecase] (uc02) at (2.0,0.25) {查询与导出};
\node[usecase] (uc03) at (2.0,-0.85) {统计分析};
\node[aicase] (uc07) at (2.0,-1.95) {决策报告};
\node[verifycase] (uc10) at (2.0,-3.7) {测试与验收};

\node[workcase] (uc05) at (7.6,0.65) {三维楼层};
\node[riskcase] (uc04) at (7.6,-0.55) {异常解释};
\node[workcase] (uc06) at (7.6,-1.75) {工单闭环};
\node[aicase] (uc08) at (7.6,-3.0) {智能问答};
\node[aicase] (uc09) at (7.6,-4.15) {MCP 调用};

\draw[-{Stealth[length=2mm]}, thick] (energy.east) -- (-0.25,0.75) |- (uc01.west);
\draw[-{Stealth[length=2mm]}, thick] (energy.east) -- (-0.25,0.75) |- (uc02.west);
\draw[-{Stealth[length=2mm]}, thick] (energy.east) -- (-0.25,0.75) |- (uc03.west);
\draw[-{Stealth[length=2mm]}, thick] (energy.east) -- (-0.25,0.75) |- (uc07.west);
\draw[-{Stealth[length=2mm]}, thick] (operator.west) -- (10.75,-0.8) |- (uc05.east);
\draw[-{Stealth[length=2mm]}, thick] (operator.west) -- (10.75,-0.8) |- (uc04.east);
\draw[-{Stealth[length=2mm]}, thick] (operator.west) -- (10.75,-0.8) |- (uc06.east);
\draw[-{Stealth[length=2mm]}, thick] (reviewer.east) -- (-0.25,-3.7) |- (uc10.west);
\draw[-{Stealth[length=2mm]}, thick] (agent.west) -- (10.75,-3.3) |- (uc08.east);
\draw[-{Stealth[length=2mm]}, thick] (agent.west) -- (10.75,-3.3) |- (uc09.east);

\draw[-{Stealth[length=2mm]}, dashed] (uc04.south) -- node[right, font=\tiny]{支撑} (uc06.north);
\node[font=\tiny, align=center, text=gray!70] at (5.25,-4.72)
{说明：三维楼层联动统计分析；异常解释支撑工单闭环；智能问答可调用 MCP 工具。};
\end{tikzpicture}
}
\caption{系统用例图}
\end{figure}

\subsection{关键用例详细说明}
\begin{longtable}{L{1.5cm}L{2.6cm}L{3.4cm}L{3.4cm}L{3.2cm}}
\caption{关键用例规约}\\
\toprule
编号 & 触发条件 & 主成功场景 & 异常或替代场景 & 验收标准 \\
\midrule
\endfirsthead
\toprule
编号 & 触发条件 & 主成功场景 & 异常或替代场景 & 验收标准 \\
\midrule
\endhead
UC-01 & 用户进入总览页 & 系统加载总览、建筑、楼层风险和趋势数据 & 后端未启动时显示接口错误提示 & KPI 与接口返回一致 \\
UC-02 & 用户设置查询条件 & 系统按建筑、楼层、时间和数量返回记录 & 条件为空时返回默认范围；无数据时显示空状态 & 表格字段完整且数量正确 \\
UC-03 & 用户点击导出 & 后端按筛选条件生成 CSV 下载 & 无数据时导出空结果或提示无记录 & CSV 可打开且字段一致 \\
UC-04 & 用户切换统计范围 & 系统刷新异常、趋势、楼层、设备和建议模块 & 健康楼层只显示健康结论 & 异常条数与筛选范围一致 \\
UC-05 & 用户点击三维楼层 & 系统跳转统计分析并锁定建筑楼层 & 点击健康楼层时不展示空图表 & 页面联动准确 \\
UC-06 & 用户点击异常解释 & 系统返回触发规则、指标证据和建议动作 & 记录不存在时返回空对象 & 解释内容可支持运维判断 \\
UC-07 & 用户创建工单 & 系统选择异常、负责人并生成处理中工单 & 重复创建时由用户确认是否继续 & 工单刷新后仍存在 \\
UC-08 & 用户完成工单 & 系统更新状态、排序和颜色 & 工单不存在时返回 404 & 已完成工单置后且绿色显示 \\
UC-09 & 用户发起问答 & 系统检索知识库并可选调用外部模型 & 外部模型失败时回退本地回答 & 回答包含项目数据或知识库依据 \\
UC-10 & AI 客户端调用 MCP & MCP Server 调用同一服务层返回数据或报告 & transport 配置错误时启动失败并提示 & MCP 工具可列出并调用 \\
\bottomrule
\end{longtable}

\section{功能需求}
\subsection{需求优先级定义}
本项目采用 Must、Should、Could 三档优先级。Must 表示最终演示和验收必须具备；Should 表示对系统完整性有明显帮助，当前版本已实现或基本实现；Could 表示可作为后续扩展。

\begin{longtable}{p{2cm}p{3.5cm}p{1.7cm}p{5.3cm}p{2.2cm}}
\caption{功能需求列表}\\
\toprule
编号 & 功能名称 & 优先级 & 需求描述 & 实现证据 \\
\midrule
\endfirsthead
\toprule
编号 & 功能名称 & 优先级 & 需求描述 & 实现证据 \\
\midrule
\endhead
FR-01 & 数据集与数据字典 & Must & 系统应提供 2,976 条建筑能耗样例数据，字段含义、类型、单位和用途应可追溯。 & \code{data/} \\
FR-02 & 数据元信息 & Must & 系统应返回字段列表、建筑清单、记录数量、建筑数量和时间范围。 & \code{/dataset-meta} \\
FR-03 & 数据总览 & Must & 系统应展示总记录数、建筑数、平均 COP、异常记录数、时间范围和能耗合计。 & \code{/overview} \\
FR-04 & 数据浏览 & Must & 系统应支持按建筑、楼层、开始时间、结束时间和数量上限筛选能耗记录。 & \code{/records} \\
FR-05 & CSV 导出 & Should & 系统应支持按当前筛选条件导出 CSV 文件。 & \code{/export/csv} \\
FR-06 & 时间序列统计 & Must & 系统应支持小时、日、周、月粒度汇总电耗、水耗、空调电耗、制冷量和 COP。 & \code{/analytics/time-summary} \\
FR-07 & 建筑对比 & Must & 系统应比较各建筑总电耗、水耗、空调电耗、制冷量和平均 COP。 & \code{/analytics/building-comparison} \\
FR-08 & COP 排名 & Must & 系统应按建筑平均 COP 排名，识别制冷效率差异。 & \code{/analytics/cop-ranking} \\
FR-09 & 异常识别 & Must & 系统应依据设备状态、建筑基线、COP 阈值和夜间负荷规则识别异常。 & \code{analysis_service.py} \\
FR-10 & 异常明细 & Must & 系统应在选定建筑/楼层/时间范围内优先展示异常记录。 & \code{/analytics/anomalies} \\
FR-11 & 异常原因统计 & Should & 系统应按异常原因统计数量，用于图表展示和诊断。 & \code{/analytics/anomaly-reasons} \\
FR-12 & 异常解释 & Must & 系统应对单条异常输出触发规则、指标证据、严重程度和建议动作。 & \code{/anomaly-explanations} \\
FR-13 & 楼层分析 & Must & 系统应派生楼层标签并统计楼层能耗、COP、异常数量和健康状态。 & \code{/floor-summary} \\
FR-14 & 设备分析 & Should & 系统应统计设备点位、设备类型、状态、最近出现时间和维护提示。 & \code{/equipment-summary} \\
FR-15 & 三维风险视图 & Must & 系统应以可拖拽三维楼宇展示楼层健康状态，并支持点击联动统计分析。 & 前端总览页 \\
FR-16 & 工单闭环 & Must & 系统应支持从异常创建工单、分配负责人、完成工单和持久化状态。 & \code{/work-orders} \\
FR-17 & 决策报告 & Should & 系统应生成运营日报，汇总风险、工单、优化建议和收益模拟。 & \code{/operation-report} \\
FR-18 & 智能问答 & Should & 系统应基于本地知识库回答问题，并在配置外部模型时增强回答。 & \code{/assistant/query} \\
FR-19 & 外部模型选择 & Should & 页面应展示可用模型选项，允许选择不同 OpenAI-compatible 提供商。 & \code{/assistant/providers} \\
FR-20 & MCP Server & Must & 系统应通过 MCP Tools/Resources 暴露数据、分析、异常解释、知识检索和问答能力。 & \code{mcp_server.py} \\
\bottomrule
\end{longtable}

\section{数据需求}
\subsection{数据规模与时间范围}
最终数据集位于 \code{data/samples/energy_records.csv}，包含 2,976 条原始记录、16 个原始字段、4 栋建筑，时间范围为 2026-03-01 00:00:00 至 2026-06-01 21:00:00。分析服务会在运行时派生楼层、区域、设备类型、COP、动态阈值、异常标记和异常原因等字段，分析层共形成 30 个字段。

\begin{table}[H]
\centering
\caption{数据集关键指标}
\begin{tabularx}{\textwidth}{p{4cm}Y}
\toprule
指标 & 当前取值 \\
\midrule
原始记录数 & 2,976 条 \\
建筑数量 & 4 栋：综合教学楼A、行政办公楼B、图书信息楼C、科研实验楼D \\
原始字段数 & 16 个 \\
分析层字段数 & 30 个 \\
楼层标签 & 1F、2F、3F、4F、5F、B1 机房、RF 屋顶 \\
建筑-楼层组合 & 14 个 \\
设备点位 & 28 个派生运维设备点位 \\
异常记录 & 548 条分析异常 \\
\bottomrule
\end{tabularx}
\end{table}

\subsection{核心字段}
\begin{longtable}{p{3.7cm}p{2.2cm}p{7.6cm}}
\caption{核心数据字段}\\
\toprule
字段 & 类型 & 说明 \\
\midrule
\endfirsthead
\toprule
字段 & 类型 & 说明 \\
\midrule
\endhead
\code{record_id} & string & 能耗记录唯一编号 \\
\code{building_id} & string & 建筑唯一编号，例如 BLD-A \\
\code{building_name} & string & 建筑展示名称 \\
\code{building_type} & string & 建筑类型，用于不同业务场景解释 \\
\code{timestamp} & datetime & 采集时间戳，当前粒度为 3 小时 \\
\code{electricity_kwh} & float & 当前时间粒度下总电耗 \\
\code{water_m3} & float & 当前时间粒度下用水量 \\
\code{hvac_kwh} & float & 空调系统电耗 \\
\code{cooling_load_kwh} & float & 制冷负荷折算值 \\
\code{environment_temp_c} & float & 室外环境温度 \\
\code{humidity_rh} & float & 相对湿度 \\
\code{occupancy_density_per_100m2} & float & 每百平方米人员密度 \\
\code{equipment_id} & string & 原始或派生设备编号 \\
\code{equipment_status} & string & 设备状态，normal 或 abnormal \\
\code{floor_label} & derived & 运行时派生楼层标签 \\
\code{average_cop} & derived & 制冷量与空调电耗比值，用于能效判断 \\
\code{is_anomaly} & derived & 是否触发异常规则 \\
\code{anomaly_reason} & derived & 异常原因说明 \\
\bottomrule
\end{longtable}

\section{外部接口需求}
\subsection{REST API}
REST API 前缀为 \code{/api/v1}，前端通过 \code{frontend/src/lib/api.js} 统一调用。接口应返回 JSON，并对无数据、参数错误或文件异常给出明确错误信息。

\begin{longtable}{p{4.2cm}p{2.2cm}p{7.5cm}}
\caption{REST API 需求}\\
\toprule
路径 & 方法 & 功能 \\
\midrule
\endfirsthead
\toprule
路径 & 方法 & 功能 \\
\midrule
\endhead
\code{/health} & GET & 健康检查 \\
\code{/overview} & GET & 系统总览指标 \\
\code{/dataset-meta} & GET & 数据集元信息 \\
\code{/buildings} & GET & 建筑下拉选项 \\
\code{/records} & GET & 能耗记录查询 \\
\code{/analytics/time-summary} & GET & 时间序列汇总 \\
\code{/analytics/building-comparison} & GET & 建筑对比 \\
\code{/analytics/cop-ranking} & GET & COP 排名 \\
\code{/analytics/anomalies} & GET & 异常明细 \\
\code{/analytics/anomaly-explanations/\{record_id\}} & GET & 单条异常解释 \\
\code{/analytics/floor-summary} & GET & 楼层汇总 \\
\code{/analytics/equipment-summary} & GET & 设备汇总 \\
\code{/analytics/operation-report} & GET & 运营日报 \\
\code{/work-orders} & GET/POST & 查询和创建持久化工单 \\
\code{/work-orders/\{id\}} & PATCH & 更新工单状态和备注 \\
\code{/assistant/query} & POST & 智能问答 \\
\code{/assistant/providers} & GET & 可用外部模型选项 \\
\bottomrule
\end{longtable}

\subsection{MCP 接口}
MCP Server 应支持 stdio 和 streamable-http 两种启动模式。MCP Tools 应覆盖数据元信息、建筑清单、能耗记录查询、总览、时间汇总、建筑对比、COP 排名、异常、异常解释、楼层汇总、设备汇总、工单建议、优化建议、运营报告、知识库检索和问答。MCP Resources 应至少包括数据集元信息、建筑清单、运营报告和知识库入口。

\subsection{环境变量接口}
真实 API Key 只能存放在仓库根目录 \code{.env}。提交仓库和最终压缩包时不得包含真实 \code{.env}。系统应提供 \code{.env.example} 作为配置模板，说明 \code{LLM_ENABLED}、\code{LLM_PROVIDER}、\code{LLM_BASE_URL}、\code{LLM_MODEL}、\code{LLM_API_KEY} 等变量。

\section{非功能需求}
\begin{longtable}{p{2cm}p{3.4cm}p{6.5cm}p{2.6cm}}
\caption{非功能需求}\\
\toprule
编号 & 类别 & 要求 & 验收方式 \\
\midrule
\endfirsthead
\toprule
编号 & 类别 & 要求 & 验收方式 \\
\midrule
\endhead
NFR-01 & 可运行性 & 同学拉取仓库后，应能根据 README 安装依赖、放置 \code{.env}、启动前后端并访问页面。 & README 演示 \\
NFR-02 & 可测试性 & 后端应提供 pytest 测试，前端应可构建，项目应提供一键验证脚本。 & \code{check-project.ps1} \\
NFR-03 & 可解释性 & 异常判断应说明触发规则和指标证据，避免黑箱结论。 & 异常解释页 \\
NFR-04 & 安全性 & 真实 API Key 不得提交；运行期工单数据不污染仓库。 & Git 配置 \\
NFR-05 & 可维护性 & REST API 和 MCP Server 应复用同一服务层，减少口径不一致。 & 代码结构 \\
NFR-06 & 性能 & 课程数据规模下，查询和图表刷新应在本地普通电脑上可交互完成。 & 人工测试 \\
NFR-07 & 可演示性 & 5 至 8 分钟内应能展示总览、查询、分析、3D、工单、报告、问答和文档。 & 演示脚本 \\
NFR-08 & 可扩展性 & 后续应可替换 CSV 为数据库、增加权限、接入真实设备和扩展模型。 & 设计说明 \\
\bottomrule
\end{longtable}

\section{需求追踪矩阵}
\begin{longtable}{p{2.2cm}p{4.2cm}p{4.2cm}p{3.5cm}}
\caption{需求追踪矩阵}\\
\toprule
需求编号 & 设计/实现位置 & 测试或验收证据 & 用户价值 \\
\midrule
\endfirsthead
\toprule
需求编号 & 设计/实现位置 & 测试或验收证据 & 用户价值 \\
\midrule
\endhead
FR-01 至 FR-04 & \code{data_loader.py}、\code{data.py}、数据浏览页 & 数据接口测试、页面查询 & 数据可查、可追溯 \\
FR-05 至 FR-08 & \code{export.py}、\code{analysis_service.py}、图表组件 & 导出测试、统计接口测试 & 管理者掌握能耗结构 \\
FR-09 至 FR-14 & 异常规则、楼层和设备分析服务 & 异常接口测试、人工验收 & 异常可解释、可定位 \\
FR-15 & \code{BuildingRiskScene.vue}、总览页联动逻辑 & 浏览器操作验证 & 风险态势直观展示 \\
FR-16 & \code{work_orders.py}、\code{work_order_store.py} & 工单接口测试、刷新验证 & 运维形成闭环 \\
FR-17 & \code{build_operation_report}、决策报告页 & 报告接口测试 & 汇报和决策支持 \\
FR-18 至 FR-19 & \code{assistant_service.py}、\code{llm_client.py} & 问答接口测试、模型配置测试 & 智能问答与模型扩展 \\
FR-20 & \code{mcp_server.py}、MCP 测试 & MCP 集成测试 & 支持 AI 客户端接入 \\
\bottomrule
\end{longtable}

\section{验收标准}
系统最终验收应满足以下条件：
\begin{enumerate}
    \item 后端自动化测试全部通过，当前基线为 75 passed。
    \item 前端生产构建成功，页面能够访问并完成核心演示路径。
    \item \code{scripts/check-project.ps1} 能正确执行并在失败时返回失败。
    \item README 能指导运行人员完成从拉取、配置、启动到访问的完整流程。
    \item SRS、SDD、SEE、SEM、测试验收、用户部署、数据接口和提交说明均已形成 PDF。
    \item 最终提交材料不包含真实 \code{.env}、API Key、\code{node_modules}、前端构建缓存和运行期工单状态。
    \item 系统演示能够覆盖总览、数据浏览、统计分析、三维楼层、工单中心、决策报告、智能问答和 MCP 说明。
\end{enumerate}

\clearpage
\section{补充需求规格}
\subsection{功能需求细化}
本节对前述功能需求进一步展开，强调每项需求的输入、处理、输出和验收口径。这样做的目的不是重复列功能，而是把课堂演示中的“能看到页面”转化为可检查、可追踪、可测试的软件需求。
\begin{longtable}{L{1.7cm}L{2.8cm}L{6.2cm}L{4.5cm}}
\caption{功能需求细化与验收口径}\\
\toprule
编号 & 功能 & 详细要求 & 验收口径 \\
\midrule
\endfirsthead
\toprule
编号 & 功能 & 详细要求 & 验收口径 \\
\midrule
\endhead
FR-01 & 数据集与字典 & 系统启动时读取样例 CSV，校验字段、时间、数值和建筑标识，数据字典应解释每个字段的含义、单位、来源和用途。 & 打开数据浏览页，字段显示完整，元信息接口返回记录数、建筑数和时间范围。 \\
FR-02 & 数据元信息 & 系统应提供建筑清单、字段清单、时间范围和数据规模，便于前端初始化筛选器并向教师说明数据来源。 & 接口返回值与 CSV 内容一致，空数据或路径错误时给出明确提示。 \\
FR-03 & 系统总览 & 总览应展示记录数、建筑数、平均 COP、异常数、总能耗、三维楼宇状态和全局趋势摘要。 & 刷新页面后总览指标与后端计算口径一致。 \\
FR-04 & 数据浏览 & 用户可按建筑、楼层、开始时间、结束时间和数量限制查询能耗记录，数据浏览只承担全量查询职责，不被 3D 点击联动污染。 & 筛选后的记录数量、建筑、楼层和时间范围均正确。 \\
FR-05 & CSV 导出 & 系统应按照当前筛选条件导出 CSV，导出字段保持与数据浏览一致，便于后续报告或表格软件分析。 & 导出文件可打开，表头和筛选结果一致。 \\
FR-06 & 时间序列统计 & 系统支持小时、日、周、月粒度汇总电耗、水耗、空调电耗、制冷量和 COP，用于趋势分析和报告。 & 切换粒度后图表刷新，空范围返回空数组而非前端错误。 \\
FR-07 & 建筑对比 & 系统按建筑汇总能耗、水耗、空调电耗、制冷量、平均 COP 和异常数量，用于横向比较。 & 不同建筑的排名和总览指标可交叉验证。 \\
FR-08 & COP 排名 & 系统按建筑平均 COP 排名，识别制冷效率差异并提示低效建筑。 & COP 数值由制冷量和空调电耗计算，排序稳定。 \\
FR-09 & 异常识别 & 系统基于设备状态、建筑动态阈值、COP 阈值和夜间负荷规则识别异常，结论应可解释。 & 异常明细中的原因能追溯到具体规则。 \\
FR-10 & 异常明细 & 统计分析页在选定建筑、楼层和时间范围后优先展示异常记录，健康楼层只显示健康结论。 & 异常条数与筛选器显示一致，健康楼层不出现无意义空图。 \\
FR-11 & 异常原因统计 & 系统按原因聚合异常数量，用于图表和报告中的风险结构说明。 & 原因统计合计数等于异常明细数。 \\
FR-12 & 异常解释 & 每条异常应给出触发规则、关键证据、严重程度、原因、结论和建议动作。 & 点击异常解释后内容足够支持运维判断。 \\
FR-13 & 楼层分析 & 系统派生楼层标签，统计楼层能耗、COP、异常数量、风险状态和设备信息。 & 3D 楼层点击可跳转到相同建筑楼层的统计分析。 \\
FR-14 & 设备分析 & 系统统计设备点位、设备类型、设备状态、最近出现时间和维护提示。 & 设备明细与异常记录可相互验证。 \\
FR-15 & 三维风险视图 & 总览页展示可旋转拖动的三维楼宇，异常楼层红色，健康楼层绿色，点击后进入统计分析。 & 拖拽、旋转和点击联动可演示。 \\
FR-16 & 工单闭环 & 用户可从异常创建工单，分配负责人，工单进入处理中，完成后变为已完成并持久化。 & 刷新后工单状态不丢失，处理中置顶，已完成置后。 \\
FR-17 & 决策报告 & 系统生成运营日报，汇总风险、工单状态、节能建议、收益模拟和优先级。 & 报告内容可用于系统展示和管理说明。 \\
FR-18 & 智能问答 & 系统基于本地知识库回答项目问题，可选调用外部模型增强，外部失败时回退本地回答。 & 无 API Key 时仍能回答基础问题。 \\
FR-19 & 模型选择 & 页面展示可用模型选项，用户选择后后端按统一 OpenAI-compatible 契约调用。 & 提供商不可用时前端显示回退结果或错误提示。 \\
FR-20 & MCP Server & MCP Server 暴露数据、统计、异常解释、报告和知识库检索能力，供 AI 客户端调用。 & MCP 工具可列出并返回与 REST API 一致的数据。 \\
\bottomrule
\end{longtable}

\subsection{典型业务场景}
系统的主要价值体现在跨页面业务闭环，而不是单个接口返回数据。以下场景覆盖最终演示和验收时最容易被教师询问的路径。
\begin{longtable}{L{1.6cm}L{3.0cm}L{6.2cm}L{4.4cm}}
\caption{典型业务场景与边界条件}\\
\toprule
编号 & 场景 & 业务说明 & 预期结果 \\
\midrule
\endfirsthead
\toprule
编号 & 场景 & 业务说明 & 预期结果 \\
\midrule
\endhead
SC-01 & 首次进入总览 & 用户打开前端后系统应自动加载总览、建筑清单、三维楼宇和关键 KPI，若后端未启动则显示接口错误提示。 & 页面不白屏，错误信息可读。 \\
SC-02 & 查询单栋建筑 & 用户选择综合教学楼 A 后，数据浏览、统计分析和楼层筛选应只使用该建筑范围。 & 表格中不出现其他建筑。 \\
SC-03 & 锁定楼层异常 & 用户在统计分析中选择 RF 屋顶，应优先显示该楼层异常明细和解释入口。 & 异常条数与楼层下拉展示一致。 \\
SC-04 & 点击健康楼层 & 用户从 3D 图点击绿色楼层，统计分析只展示健康提示，不展示趋势、异常和推荐模块。 & 页面说明该楼层健康且无空图表。 \\
SC-05 & 点击异常楼层 & 用户从 3D 图点击红色楼层，系统跳转统计分析并自动填入建筑和楼层。 & 筛选标签显示目标建筑和楼层。 \\
SC-06 & 创建工单 & 用户选择某条异常并分配给管理员，系统生成处理中工单并置顶展示。 & 工单编号、负责人、异常引用均存在。 \\
SC-07 & 完成工单 & 用户点击完成按钮后，工单状态变为已完成，颜色变绿并移动到列表后部。 & 刷新页面后状态仍为已完成。 \\
SC-08 & 导出数据 & 用户在数据浏览页完成筛选后导出 CSV，系统按当前条件生成文件。 & 文件可打开，记录范围正确。 \\
SC-09 & 生成日报 & 管理人员进入决策报告页，系统按当前数据生成运营日报和节能建议。 & 报告包含风险摘要和优先事项。 \\
SC-10 & 本地问答 & 无外部模型密钥时，用户询问异常规则或演示路径，系统使用本地知识库回答。 & 回答不依赖网络且包含项目内容。 \\
SC-11 & 外部模型问答 & 配置外部模型后，用户选择模型并提问，系统把项目上下文作为提示输入。 & 回答更自然，失败时有回退。 \\
SC-12 & MCP 调用 & 支持 MCP 的客户端调用工具获取建筑清单或异常解释。 & 工具返回 JSON 或文本资源。 \\
SC-13 & 时间范围为空 & 用户选择没有记录的时间范围，系统显示空状态而不是前端错误。 & 空状态文字明确。 \\
SC-14 & 异常不存在 & 用户请求不存在的异常解释，后端返回空对象或 404，前端显示提示。 & 系统不崩溃。 \\
SC-15 & 模型密钥缺失 & 模型提供商被启用但 API Key 未配置，系统提示配置问题并回退。 & 真实密钥不出现在页面或日志。 \\
SC-16 & 数据文件缺失 & 后端启动时无法找到 CSV，系统应给出明确错误，便于定位部署问题。 & 启动日志或接口错误可追踪。 \\
SC-17 & 工单文件缺失 & 运行目录没有工单 JSON 时，系统自动创建或返回空列表。 & 不会影响正常演示。 \\
SC-18 & 重复创建工单 & 用户对同一异常重复创建工单时，系统允许创建但保留异常引用，便于演示多负责人场景。 & 工单编号不冲突。 \\
SC-19 & 接口响应慢 & 前端应显示加载状态，避免用户误以为页面无反应。 & 加载状态覆盖核心卡片。 \\
SC-20 & 最终验收 & 评审人员按演示脚本走完整路径，系统应覆盖总览、查询、分析、3D、工单、报告、问答和文档。 & 演示路径闭环完整。 \\
\bottomrule
\end{longtable}

\subsection{数据质量与安全约束}
数据质量约束用于保证分析结果可信，安全约束用于保证仓库提交和课堂演示不会泄露真实密钥。
\begin{longtable}{L{1.7cm}L{3.0cm}L{6.0cm}L{4.5cm}}
\caption{数据质量、安全与合规约束}\\
\toprule
编号 & 约束项 & 约束说明 & 处理方式 \\
\midrule
\endfirsthead
\toprule
编号 & 约束项 & 约束说明 & 处理方式 \\
\midrule
\endhead
DQ-01 & 字段完整 & CSV 必须包含建筑、时间、能耗、水耗、空调电耗、制冷量、环境温湿度、人员密度和设备状态等核心字段。 & 启动或加载时校验字段，不满足时拒绝继续分析。 \\
DQ-02 & 时间可解析 & 时间字段必须可解析为统一时间戳，前端筛选必须使用相同时间口径。 & 数据加载层统一转换。 \\
DQ-03 & 数值非空 & 电耗、水耗、空调电耗、制冷量等指标不能为空，否则会影响 COP 和异常计算。 & 异常值在加载阶段标记或排除。 \\
DQ-04 & 建筑标识一致 & 同一建筑的编号和名称需要稳定，避免总览、统计和 3D 展示口径不一致。 & 建筑清单由后端统一生成。 \\
DQ-05 & 楼层派生可解释 & 楼层来自设备和建筑场景推断，必须能说明为什么冷水机组映射到 B1 或屋顶设备映射到 RF。 & SDD 中记录派生规则。 \\
DQ-06 & 异常原因唯一主因 & 一条记录可触发多条规则，但页面优先展示一个主因并保留证据说明。 & 分析服务按优先级排序。 \\
DQ-07 & 工单持久化 & 工单状态必须写入运行目录，刷新后不丢失，但不进入 Git 仓库。 & \code{data/runtime/} 加入忽略。 \\
DQ-08 & 密钥不入库 & 真实 API Key 只能放在本地 \code{.env}，不得写入源码、文档、PDF 或提交材料。 & 发布前执行模式搜索。 \\
DQ-09 & 外部模型可关闭 & 没有外部模型时系统仍应使用本地知识库回答，避免演示依赖网络。 & 通过 \code{LLM_ENABLED} 控制。 \\
DQ-10 & 文档正式口吻 & 最终文档应面向课程评审，不应出现对话式说明或个人任务口吻。 & 生成正式 PDF 前统一确认。 \\
\bottomrule
\end{longtable}

\subsection{需求可追踪性说明}
需求可追踪性用于说明需求文档是否能够支撑设计、开发、测试和最终验收。每一项都要求能在系统、代码或文档中找到证据。
\begin{longtable}{L{1.7cm}L{3.0cm}L{6.2cm}L{4.4cm}}
\caption{需求可追踪性说明}\\
\toprule
编号 & 需求项 & 追踪说明 & 验收依据 \\
\midrule
\endfirsthead
\toprule
编号 & 需求项 & 追踪说明 & 验收依据 \\
\midrule
\endhead
SR-01 & 项目目标 & “项目目标”已经在需求、界面、接口或验收说明中被明确描述，避免只停留在口头说明。 & 能够在正式文档或系统页面中找到“项目目标”对应证据。 \\
SR-02 & 用户角色 & “用户角色”已经在需求、界面、接口或验收说明中被明确描述，避免只停留在口头说明。 & 能够在正式文档或系统页面中找到“用户角色”对应证据。 \\
SR-03 & 业务边界 & “业务边界”已经在需求、界面、接口或验收说明中被明确描述，避免只停留在口头说明。 & 能够在正式文档或系统页面中找到“业务边界”对应证据。 \\
SR-04 & 数据来源 & “数据来源”已经在需求、界面、接口或验收说明中被明确描述，避免只停留在口头说明。 & 能够在正式文档或系统页面中找到“数据来源”对应证据。 \\
SR-05 & 字段含义 & “字段含义”已经在需求、界面、接口或验收说明中被明确描述，避免只停留在口头说明。 & 能够在正式文档或系统页面中找到“字段含义”对应证据。 \\
SR-06 & 建筑范围 & “建筑范围”已经在需求、界面、接口或验收说明中被明确描述，避免只停留在口头说明。 & 能够在正式文档或系统页面中找到“建筑范围”对应证据。 \\
SR-07 & 楼层范围 & “楼层范围”已经在需求、界面、接口或验收说明中被明确描述，避免只停留在口头说明。 & 能够在正式文档或系统页面中找到“楼层范围”对应证据。 \\
SR-08 & 时间范围 & “时间范围”已经在需求、界面、接口或验收说明中被明确描述，避免只停留在口头说明。 & 能够在正式文档或系统页面中找到“时间范围”对应证据。 \\
SR-09 & 查询条件 & “查询条件”已经在需求、界面、接口或验收说明中被明确描述，避免只停留在口头说明。 & 能够在正式文档或系统页面中找到“查询条件”对应证据。 \\
SR-10 & 导出功能 & “导出功能”已经在需求、界面、接口或验收说明中被明确描述，避免只停留在口头说明。 & 能够在正式文档或系统页面中找到“导出功能”对应证据。 \\
SR-11 & 总览指标 & “总览指标”已经在需求、界面、接口或验收说明中被明确描述，避免只停留在口头说明。 & 能够在正式文档或系统页面中找到“总览指标”对应证据。 \\
SR-12 & 趋势统计 & “趋势统计”已经在需求、界面、接口或验收说明中被明确描述，避免只停留在口头说明。 & 能够在正式文档或系统页面中找到“趋势统计”对应证据。 \\
SR-13 & 建筑对比 & “建筑对比”已经在需求、界面、接口或验收说明中被明确描述，避免只停留在口头说明。 & 能够在正式文档或系统页面中找到“建筑对比”对应证据。 \\
SR-14 & COP 排名 & “COP 排名”已经在需求、界面、接口或验收说明中被明确描述，避免只停留在口头说明。 & 能够在正式文档或系统页面中找到“COP 排名”对应证据。 \\
SR-15 & 异常规则 & “异常规则”已经在需求、界面、接口或验收说明中被明确描述，避免只停留在口头说明。 & 能够在正式文档或系统页面中找到“异常规则”对应证据。 \\
SR-16 & 异常明细 & “异常明细”已经在需求、界面、接口或验收说明中被明确描述，避免只停留在口头说明。 & 能够在正式文档或系统页面中找到“异常明细”对应证据。 \\
SR-17 & 异常解释 & “异常解释”已经在需求、界面、接口或验收说明中被明确描述，避免只停留在口头说明。 & 能够在正式文档或系统页面中找到“异常解释”对应证据。 \\
SR-18 & 健康楼层 & “健康楼层”已经在需求、界面、接口或验收说明中被明确描述，避免只停留在口头说明。 & 能够在正式文档或系统页面中找到“健康楼层”对应证据。 \\
SR-19 & 设备分析 & “设备分析”已经在需求、界面、接口或验收说明中被明确描述，避免只停留在口头说明。 & 能够在正式文档或系统页面中找到“设备分析”对应证据。 \\
SR-20 & 三维视图 & “三维视图”已经在需求、界面、接口或验收说明中被明确描述，避免只停留在口头说明。 & 能够在正式文档或系统页面中找到“三维视图”对应证据。 \\
SR-21 & 楼层联动 & “楼层联动”已经在需求、界面、接口或验收说明中被明确描述，避免只停留在口头说明。 & 能够在正式文档或系统页面中找到“楼层联动”对应证据。 \\
SR-22 & 工单创建 & “工单创建”已经在需求、界面、接口或验收说明中被明确描述，避免只停留在口头说明。 & 能够在正式文档或系统页面中找到“工单创建”对应证据。 \\
SR-23 & 工单完成 & “工单完成”已经在需求、界面、接口或验收说明中被明确描述，避免只停留在口头说明。 & 能够在正式文档或系统页面中找到“工单完成”对应证据。 \\
SR-24 & 工单持久化 & “工单持久化”已经在需求、界面、接口或验收说明中被明确描述，避免只停留在口头说明。 & 能够在正式文档或系统页面中找到“工单持久化”对应证据。 \\
SR-25 & 运营日报 & “运营日报”已经在需求、界面、接口或验收说明中被明确描述，避免只停留在口头说明。 & 能够在正式文档或系统页面中找到“运营日报”对应证据。 \\
SR-26 & 节能建议 & “节能建议”已经在需求、界面、接口或验收说明中被明确描述，避免只停留在口头说明。 & 能够在正式文档或系统页面中找到“节能建议”对应证据。 \\
SR-27 & 收益模拟 & “收益模拟”已经在需求、界面、接口或验收说明中被明确描述，避免只停留在口头说明。 & 能够在正式文档或系统页面中找到“收益模拟”对应证据。 \\
SR-28 & 智能问答 & “智能问答”已经在需求、界面、接口或验收说明中被明确描述，避免只停留在口头说明。 & 能够在正式文档或系统页面中找到“智能问答”对应证据。 \\
SR-29 & 模型选择 & “模型选择”已经在需求、界面、接口或验收说明中被明确描述，避免只停留在口头说明。 & 能够在正式文档或系统页面中找到“模型选择”对应证据。 \\
SR-30 & 本地回退 & “本地回退”已经在需求、界面、接口或验收说明中被明确描述，避免只停留在口头说明。 & 能够在正式文档或系统页面中找到“本地回退”对应证据。 \\
SR-31 & MCP 工具 & “MCP 工具”已经在需求、界面、接口或验收说明中被明确描述，避免只停留在口头说明。 & 能够在正式文档或系统页面中找到“MCP 工具”对应证据。 \\
SR-32 & MCP 资源 & “MCP 资源”已经在需求、界面、接口或验收说明中被明确描述，避免只停留在口头说明。 & 能够在正式文档或系统页面中找到“MCP 资源”对应证据。 \\
SR-33 & 接口错误 & “接口错误”已经在需求、界面、接口或验收说明中被明确描述，避免只停留在口头说明。 & 能够在正式文档或系统页面中找到“接口错误”对应证据。 \\
SR-34 & 空数据状态 & “空数据状态”已经在需求、界面、接口或验收说明中被明确描述，避免只停留在口头说明。 & 能够在正式文档或系统页面中找到“空数据状态”对应证据。 \\
SR-35 & 权限边界 & “权限边界”已经在需求、界面、接口或验收说明中被明确描述，避免只停留在口头说明。 & 能够在正式文档或系统页面中找到“权限边界”对应证据。 \\
SR-36 & 密钥安全 & “密钥安全”已经在需求、界面、接口或验收说明中被明确描述，避免只停留在口头说明。 & 能够在正式文档或系统页面中找到“密钥安全”对应证据。 \\
SR-37 & 部署说明 & “部署说明”已经在需求、界面、接口或验收说明中被明确描述，避免只停留在口头说明。 & 能够在正式文档或系统页面中找到“部署说明”对应证据。 \\
SR-38 & 测试口径 & “测试口径”已经在需求、界面、接口或验收说明中被明确描述，避免只停留在口头说明。 & 能够在正式文档或系统页面中找到“测试口径”对应证据。 \\
SR-39 & 验收口径 & “验收口径”已经在需求、界面、接口或验收说明中被明确描述，避免只停留在口头说明。 & 能够在正式文档或系统页面中找到“验收口径”对应证据。 \\
SR-40 & 系统展示路径 & “系统展示路径”已经在需求、界面、接口或验收说明中被明确描述，避免只停留在口头说明。 & 能够在正式文档或系统页面中找到“系统展示路径”对应证据。 \\
SR-41 & 性能边界 & “性能边界”已经在需求、界面、接口或验收说明中被明确描述，避免只停留在口头说明。 & 能够在正式文档或系统页面中找到“性能边界”对应证据。 \\
SR-42 & 扩展方向 & “扩展方向”已经在需求、界面、接口或验收说明中被明确描述，避免只停留在口头说明。 & 能够在正式文档或系统页面中找到“扩展方向”对应证据。 \\
SR-43 & 维护责任 & “维护责任”已经在需求、界面、接口或验收说明中被明确描述，避免只停留在口头说明。 & 能够在正式文档或系统页面中找到“维护责任”对应证据。 \\
SR-44 & 交付材料 & “交付材料”已经在需求、界面、接口或验收说明中被明确描述，避免只停留在口头说明。 & 能够在正式文档或系统页面中找到“交付材料”对应证据。 \\
SR-45 & 文档一致性 & “文档一致性”已经在需求、界面、接口或验收说明中被明确描述，避免只停留在口头说明。 & 能够在正式文档或系统页面中找到“文档一致性”对应证据。 \\
SR-46 & 风险说明 & “风险说明”已经在需求、界面、接口或验收说明中被明确描述，避免只停留在口头说明。 & 能够在正式文档或系统页面中找到“风险说明”对应证据。 \\
SR-47 & 配置项 & “配置项”已经在需求、界面、接口或验收说明中被明确描述，避免只停留在口头说明。 & 能够在正式文档或系统页面中找到“配置项”对应证据。 \\
SR-48 & 日志提示 & “日志提示”已经在需求、界面、接口或验收说明中被明确描述，避免只停留在口头说明。 & 能够在正式文档或系统页面中找到“日志提示”对应证据。 \\
SR-49 & 最终结论 & “最终结论”已经在需求、界面、接口或验收说明中被明确描述，避免只停留在口头说明。 & 能够在正式文档或系统页面中找到“最终结论”对应证据。 \\
SR-50 & 课程验收点 & “课程验收点”已经在需求、界面、接口或验收说明中被明确描述，避免只停留在口头说明。 & 能够在正式文档或系统页面中找到“课程验收点”对应证据。 \\
\bottomrule
\end{longtable}

\subsection{需求合理性说明}
本节说明关键需求的来源、业务价值、验收方式和与系统功能的对应关系，使需求文档不仅描述“系统有什么功能”，也说明这些功能为什么纳入本期范围。
\begin{longtable}{L{1.7cm}L{4.4cm}L{6.0cm}L{3.4cm}}
\caption{需求合理性说明}\\
\toprule
编号 & 需求主题 & 合理性说明 & 关联证据 \\
\midrule
\endfirsthead
\toprule
编号 & 需求主题 & 合理性说明 & 关联证据 \\
\midrule
\endhead
RA-01 & 建筑能源管理场景 & “建筑能源管理场景”与建筑能源管理、异常处置、系统运行或后续扩展有关，能够支撑本项目的业务闭环和工程完整性。 & SRS、页面展示或接口返回。 \\
RA-02 & 楼层维度分析 & “楼层维度分析”与建筑能源管理、异常处置、系统运行或后续扩展有关，能够支撑本项目的业务闭环和工程完整性。 & SRS、页面展示或接口返回。 \\
RA-03 & 异常优先展示 & “异常优先展示”与建筑能源管理、异常处置、系统运行或后续扩展有关，能够支撑本项目的业务闭环和工程完整性。 & SRS、页面展示或接口返回。 \\
RA-04 & 健康楼层简化展示 & “健康楼层简化展示”与建筑能源管理、异常处置、系统运行或后续扩展有关，能够支撑本项目的业务闭环和工程完整性。 & SRS、页面展示或接口返回。 \\
RA-05 & 工单创建后直接处理中 & “工单创建后直接处理中”与建筑能源管理、异常处置、系统运行或后续扩展有关，能够支撑本项目的业务闭环和工程完整性。 & SRS、页面展示或接口返回。 \\
RA-06 & 工单状态持久化 & “工单状态持久化”与建筑能源管理、异常处置、系统运行或后续扩展有关，能够支撑本项目的业务闭环和工程完整性。 & SRS、页面展示或接口返回。 \\
RA-07 & 数据浏览独立查询 & “数据浏览独立查询”与建筑能源管理、异常处置、系统运行或后续扩展有关，能够支撑本项目的业务闭环和工程完整性。 & SRS、页面展示或接口返回。 \\
RA-08 & 统计分析承担异常诊断 & “统计分析承担异常诊断”与建筑能源管理、异常处置、系统运行或后续扩展有关，能够支撑本项目的业务闭环和工程完整性。 & SRS、页面展示或接口返回。 \\
RA-09 & COP 指标 & “COP 指标”与建筑能源管理、异常处置、系统运行或后续扩展有关，能够支撑本项目的业务闭环和工程完整性。 & SRS、页面展示或接口返回。 \\
RA-10 & 设备状态规则 & “设备状态规则”与建筑能源管理、异常处置、系统运行或后续扩展有关，能够支撑本项目的业务闭环和工程完整性。 & SRS、页面展示或接口返回。 \\
RA-11 & 夜间负荷规则 & “夜间负荷规则”与建筑能源管理、异常处置、系统运行或后续扩展有关，能够支撑本项目的业务闭环和工程完整性。 & SRS、页面展示或接口返回。 \\
RA-12 & 可解释规则算法 & “可解释规则算法”与建筑能源管理、异常处置、系统运行或后续扩展有关，能够支撑本项目的业务闭环和工程完整性。 & SRS、页面展示或接口返回。 \\
RA-13 & CSV 导出 & “CSV 导出”与建筑能源管理、异常处置、系统运行或后续扩展有关，能够支撑本项目的业务闭环和工程完整性。 & SRS、页面展示或接口返回。 \\
RA-14 & 运营日报 & “运营日报”与建筑能源管理、异常处置、系统运行或后续扩展有关，能够支撑本项目的业务闭环和工程完整性。 & SRS、页面展示或接口返回。 \\
RA-15 & 本地知识库 & “本地知识库”与建筑能源管理、异常处置、系统运行或后续扩展有关，能够支撑本项目的业务闭环和工程完整性。 & SRS、页面展示或接口返回。 \\
RA-16 & 外部模型可选 & “外部模型可选”与建筑能源管理、异常处置、系统运行或后续扩展有关，能够支撑本项目的业务闭环和工程完整性。 & SRS、页面展示或接口返回。 \\
RA-17 & MCP 与 REST 共用服务层 & “MCP 与 REST 共用服务层”与建筑能源管理、异常处置、系统运行或后续扩展有关，能够支撑本项目的业务闭环和工程完整性。 & SRS、页面展示或接口返回。 \\
RA-18 & 真实密钥不入库 & “真实密钥不入库”与建筑能源管理、异常处置、系统运行或后续扩展有关，能够支撑本项目的业务闭环和工程完整性。 & SRS、页面展示或接口返回。 \\
RA-19 & 多月样例数据 & “多月样例数据”与建筑能源管理、异常处置、系统运行或后续扩展有关，能够支撑本项目的业务闭环和工程完整性。 & SRS、页面展示或接口返回。 \\
RA-20 & 自动检测定位 & “自动检测定位”与建筑能源管理、异常处置、系统运行或后续扩展有关，能够支撑本项目的业务闭环和工程完整性。 & SRS、页面展示或接口返回。 \\
RA-21 & 数据与页面解耦 & “数据与页面解耦”与建筑能源管理、异常处置、系统运行或后续扩展有关，能够支撑本项目的业务闭环和工程完整性。 & SRS、页面展示或接口返回。 \\
RA-22 & 异常条数一致性 & “异常条数一致性”与建筑能源管理、异常处置、系统运行或后续扩展有关，能够支撑本项目的业务闭环和工程完整性。 & SRS、页面展示或接口返回。 \\
RA-23 & 楼层联动有效性 & “楼层联动有效性”与建筑能源管理、异常处置、系统运行或后续扩展有关，能够支撑本项目的业务闭环和工程完整性。 & SRS、页面展示或接口返回。 \\
RA-24 & 刷新后状态保留 & “刷新后状态保留”与建筑能源管理、异常处置、系统运行或后续扩展有关，能够支撑本项目的业务闭环和工程完整性。 & SRS、页面展示或接口返回。 \\
RA-25 & 模型失败回退 & “模型失败回退”与建筑能源管理、异常处置、系统运行或后续扩展有关，能够支撑本项目的业务闭环和工程完整性。 & SRS、页面展示或接口返回。 \\
RA-26 & 接口契约稳定性 & “接口契约稳定性”与建筑能源管理、异常处置、系统运行或后续扩展有关，能够支撑本项目的业务闭环和工程完整性。 & SRS、页面展示或接口返回。 \\
RA-27 & 系统可复现运行 & “系统可复现运行”与建筑能源管理、异常处置、系统运行或后续扩展有关，能够支撑本项目的业务闭环和工程完整性。 & SRS、页面展示或接口返回。 \\
RA-28 & 核心功能测试覆盖 & “核心功能测试覆盖”与建筑能源管理、异常处置、系统运行或后续扩展有关，能够支撑本项目的业务闭环和工程完整性。 & SRS、页面展示或接口返回。 \\
RA-29 & 文档与系统一致性 & “文档与系统一致性”与建筑能源管理、异常处置、系统运行或后续扩展有关，能够支撑本项目的业务闭环和工程完整性。 & SRS、页面展示或接口返回。 \\
RA-30 & 经济分析关联性 & “经济分析关联性”与建筑能源管理、异常处置、系统运行或后续扩展有关，能够支撑本项目的业务闭环和工程完整性。 & SRS、页面展示或接口返回。 \\
RA-31 & README 运行说明 & “README 运行说明”与建筑能源管理、异常处置、系统运行或后续扩展有关，能够支撑本项目的业务闭环和工程完整性。 & SRS、页面展示或接口返回。 \\
RA-32 & SRS 需求基线 & “SRS 需求基线”与建筑能源管理、异常处置、系统运行或后续扩展有关，能够支撑本项目的业务闭环和工程完整性。 & SRS、页面展示或接口返回。 \\
RA-33 & SDD 设计基线 & “SDD 设计基线”与建筑能源管理、异常处置、系统运行或后续扩展有关，能够支撑本项目的业务闭环和工程完整性。 & SRS、页面展示或接口返回。 \\
RA-34 & SEE 经济评价 & “SEE 经济评价”与建筑能源管理、异常处置、系统运行或后续扩展有关，能够支撑本项目的业务闭环和工程完整性。 & SRS、页面展示或接口返回。 \\
RA-35 & SEM 工程管理 & “SEM 工程管理”与建筑能源管理、异常处置、系统运行或后续扩展有关，能够支撑本项目的业务闭环和工程完整性。 & SRS、页面展示或接口返回。 \\
RA-36 & 系统启动路径 & “系统启动路径”与建筑能源管理、异常处置、系统运行或后续扩展有关，能够支撑本项目的业务闭环和工程完整性。 & SRS、页面展示或接口返回。 \\
RA-37 & 页面刷新行为 & “页面刷新行为”与建筑能源管理、异常处置、系统运行或后续扩展有关，能够支撑本项目的业务闭环和工程完整性。 & SRS、页面展示或接口返回。 \\
RA-38 & 源码结构 & “源码结构”与建筑能源管理、异常处置、系统运行或后续扩展有关，能够支撑本项目的业务闭环和工程完整性。 & SRS、页面展示或接口返回。 \\
RA-39 & API 契约 & “API 契约”与建筑能源管理、异常处置、系统运行或后续扩展有关，能够支撑本项目的业务闭环和工程完整性。 & SRS、页面展示或接口返回。 \\
RA-40 & 交付材料结构 & “交付材料结构”与建筑能源管理、异常处置、系统运行或后续扩展有关，能够支撑本项目的业务闭环和工程完整性。 & SRS、页面展示或接口返回。 \\
RA-41 & 真实楼宇数据扩展 & “真实楼宇数据扩展”与建筑能源管理、异常处置、系统运行或后续扩展有关，能够支撑本项目的业务闭环和工程完整性。 & SRS、页面展示或接口返回。 \\
RA-42 & 用户权限扩展 & “用户权限扩展”与建筑能源管理、异常处置、系统运行或后续扩展有关，能够支撑本项目的业务闭环和工程完整性。 & SRS、页面展示或接口返回。 \\
RA-43 & 数据库替换扩展 & “数据库替换扩展”与建筑能源管理、异常处置、系统运行或后续扩展有关，能够支撑本项目的业务闭环和工程完整性。 & SRS、页面展示或接口返回。 \\
RA-44 & 预测模型扩展 & “预测模型扩展”与建筑能源管理、异常处置、系统运行或后续扩展有关，能够支撑本项目的业务闭环和工程完整性。 & SRS、页面展示或接口返回。 \\
RA-45 & 真实工单系统扩展 & “真实工单系统扩展”与建筑能源管理、异常处置、系统运行或后续扩展有关，能够支撑本项目的业务闭环和工程完整性。 & SRS、页面展示或接口返回。 \\
\bottomrule
\end{longtable}

\subsection{需求变更与版本控制说明}
项目在最终打磨阶段根据系统展示效果、课程要求和系统完整性进行了多轮需求修正。以下记录说明主要变更的原因和需求层影响。
\begin{longtable}{L{1.7cm}L{4.2cm}L{6.6cm}L{3.0cm}}
\caption{需求变更与版本控制说明}\\
\toprule
编号 & 变更项 & 变更原因与影响 & 状态 \\
\midrule
\endfirsthead
\toprule
编号 & 变更项 & 变更原因与影响 & 状态 \\
\midrule
\endhead
C-01 & 从单日数据扩展为多月数据 & 解决数据量太少、趋势不足的问题，使系统能展示 2026-03-01 至 2026-06-01 的连续记录。 & 已纳入最终版 \\
C-02 & 增加楼层筛选 & 把建筑级分析细化到楼层级，支撑 3D 楼层点击和异常定位。 & 已纳入最终版 \\
C-03 & 统计分析承接异常查询 & 将异常相关内容集中在统计分析页，数据浏览保持全量查询职责。 & 已纳入最终版 \\
C-04 & 健康楼层简化展示 & 健康楼层只显示健康提示，避免空图表降低页面表达质量。 & 已纳入最终版 \\
C-05 & 增加三维楼宇视图 & 通过立体楼宇展示楼层风险，使总览页更直观。 & 已纳入最终版 \\
C-06 & 增加工单持久化 & 避免刷新后工单消失，提高业务闭环可信度。 & 已纳入最终版 \\
C-07 & 调整工单状态 & 去掉无意义的待处理状态，创建后直接进入处理中。 & 已纳入最终版 \\
C-08 & 增加决策报告 & 把异常、工单和收益建议汇总为管理视角输出。 & 已纳入最终版 \\
C-09 & 增加外部模型选择 & 支持多个 OpenAI-compatible 模型提供商，但保留本地回退。 & 已纳入最终版 \\
C-10 & 增加 MCP Server & 让 AI 客户端可直接调用项目数据和分析能力。 & 已纳入最终版 \\
C-11 & 增加正式 LaTeX 文档 & 把 Markdown 过程资料升级为可提交的正式 PDF。 & 已纳入最终版 \\
C-12 & 重画关键流程图 & 解决图形重叠和箭头混乱问题，提高文档可读性。 & 已纳入最终版 \\
C-13 & 收紧表格版式 & 避免长表超出页面宽度。 & 已纳入最终版 \\
C-14 & 增加交付材料目录 & 将源码、正式 PDF、LaTeX 源和展示材料分区整理。 & 已纳入最终版 \\
C-15 & 增加安全控制 & 确保交付材料不包含真实 .env 和 API Key。 & 已纳入最终版 \\
\bottomrule
\end{longtable}

\end{document}
